
classe de \(f^*M_K\) est dans \(\mathcal L\), les trois coefficients \(a_r,a_{r-1},a_{r-2}\) restent born\'es gr\^ace \`a 1.13.

Soit \(H={O}_Y(m)\) un multiple de \({O}_Y(1)\) qui est tr\`es ample. Alors dans les polyn\^omes de Hilbert \(\chi(M_K(mn))=\sum_{i=0}^{r}a_i^{(m)}\binom{n+i}{i}\) relativement \`a \(H\) des \(M_K\) dans \((f^*)^{-1}(\mathcal L)\), les trois coefficients \(a_r^{(m)},a_{r-1}^{(m)},a_{r-2}^{(m)}\) restent born\'es gr\^ace \`a 2.10. Par cons\'equent \((f^*)^{-1}(\mathcal L)\) est limit\'ee en vertu de 2.11\,; d'o\`u la premi\`ere partie de 3.8.

Quant \`a la deuxi\`eme partie de 3.8, la repr\'esentabilit\'e ``g\'en\'erique'' des foncteurs de Picard a \'et\'e \'etablie dans XII 1.2, ainsi que leur s\'eparation. Par un raisonnement similaire \`a celui de XII 2, utilisant 3.9 et XII 1.1, on se ram\`ene au cas \(Y\) normal sur \(S\). Or, pour que le morphisme \(\underline{\Pic}_{Y/S}\to\underline{\Pic}_{X/S}\) soit quasi-affine, il suffit qu'il soit (s\'epar\'e et) quasi-fini (EGA IV 8.12.8), ce qui se v\'erifie sur les fibres g\'eom\'etriques, et r\'esulte donc du lemme suivant~:

\textbf{Lemme 3.11.} \emph{Soient \(k\) un corps alg\'ebriquement clos, \(Y\) un \(k\)-sch\'ema normal, projectif muni d'un faisceau ample \({O}_Y(1)\), et \(X\) le sch\'ema des z\'eros d'une section \(\varphi\) de \({O}_Y(1)\). On suppose que \(Y\) soit de dimension \(\ge3\) (resp. \(\ge2\)). Alors le morphisme \(\Pic_{Y/k}\longrightarrow\Pic_{X/k}\) (resp. le morphisme \(\Pic^\tau_{Y/k}\to\Pic^\tau_{X/k}\) (voir 4.1)) est de type fini et a un noyau \(N\) qui est un groupe fini (non-r\'eduit en g\'en\'eral).}

En effet, on peut \'evidemment supposer \(\varphi\ne0\)\,; donc \(\varphi\) est r\'eguli\`ere, \(Y\) \'etant int\`egre. Le morphisme envisag\'e est de type fini gr\^ace \`a la premi\`ere partie de 3.8 (resp. gr\^ace \`a 4.7 (iii)). Par suite le noyau \(N\) est un groupe alg\'ebrique, donc contenu dans \(\Pic^\tau_{Y/k}\), donc est le m\^eme
